% Abhakseite „Das kann ich“ – Zone nur im Gesamt (\mitzone)
%% TODO Ich-kann-Sätze: Platzhalter wie in den Aufgabendateien
\begin{abhakseite}
\ifdefined\mitzone
\abhakgruppe{Kennst du schon}
\abhak{1}{Ebenengleichungen lesen (Koordinatenform, Normalenvektor ablesen, Parameterform)}
\abhak{2}{Geradengleichungen lesen und den allgemeinen Geradenpunkt bilden (Stützpunkt plus Parameter mal Richtungsvektor)}
\abhak{3}{Skalarprodukt berechnen und Orthogonalität erkennen}
\abhak{4}{Gleichungen lösen}
\abhak{5}{Terme mit einem Parameter umformen und Fallunterscheidungen führen}
\abhak{6}{Punkte im räumlichen Koordinatensystem lesen und Bereiche deuten (Koordinatenschranken, Wandmaße)}
\abhak{7}{Ebenengleichungen lesen (Koordinatenform, Normalenvektor ablesen, Parameterform) – Fehler finden}
\abhak{8}{Ebenengleichungen lesen (Koordinatenform, Normalenvektor ablesen, Parameterform) – selbst rechnen}
\fi
\abhakgruppe{Einheit 1 · Punktprobe und Seitenlage}
\abhak{9}{Wer und wogegen?}
\abhak{10}{Punktprobe und Seitenlage}
\abhak{11}{Punktprobe und Seitenlage – Fehler finden, Begründen}
\abhakgruppe{Einheit 2 · Parameter aus der Lagebedingung}
\abhak{12}{Parameter aus der Lagebedingung}
\abhak{13}{Parameter aus der Lagebedingung – Fehler finden, Begründen}
\abhakgruppe{Einheit 3 · Gerade und Ebene}
\abhak{14}{Gerade und Ebene}
\abhak{15}{Gerade und Ebene – Fehler finden, Begründen}
\abhakgruppe{Einheit 4 · Lagebefunde im Sachzusammenhang}
\abhak{16}{Lagebefunde im Sachzusammenhang}
\abhak{17}{Gerade und Ebene: Länge eines Schattens auf einer Dachfläche über Schnitt von Lichtgerade und Ebene beschreiben}
\abhak{18}{Lagebefunde im Sachzusammenhang – Fehler finden, Begründen}
\end{abhakseite}
